\documentclass[a4paper,11pt]{ltjsarticle}
\usepackage{luatexja}
\usepackage{luatexja-fontspec}
\usepackage{amsmath,amssymb,amsthm}
\usepackage{mathtools}
\usepackage{booktabs}
\usepackage{longtable}
\usepackage{array}
\usepackage{hyperref}
\usepackage{graphicx}
\usepackage{float}
\usepackage{geometry}
\geometry{margin=25mm}

\hypersetup{
  colorlinks=true,
  linkcolor=blue,
  urlcolor=blue,
  citecolor=blue
}

\theoremstyle{definition}
\newtheorem{theorem}{定理}[section]
\newtheorem{proposition}[theorem]{命題}
\newtheorem{lemma}[theorem]{補題}
\newtheorem{corollary}[theorem]{系}
\newtheorem{definition}[theorem]{定義}
\newtheorem{remark}[theorem]{注意}
\newtheorem{observation}[theorem]{観察}

\setlength{\parindent}{1em}
\setlength{\parskip}{0.5em}

\providecommand{\tightlist}{\setlength{\itemsep}{0pt}\setlength{\parskip}{0pt}}
\begin{document}

\section{ゼロ次元から四次元主観空間の考察}\label{ux30bcux30edux6b21ux5143ux304bux3089ux56dbux6b21ux5143ux4e3bux89b3ux7a7aux9593ux306eux8003ux5bdf}

\begin{center}\rule{0.5\linewidth}{0.5pt}\end{center}

\textbf{著者}: 木原 範昭 (WF System Co., Ltd.)

\textbf{日付}: 2026年4月

\textbf{DOI：} \href{https://doi.org/10.5281/zenodo.19533292}{10.5281/zenodo.19533292}

\begin{center}\rule{0.5\linewidth}{0.5pt}\end{center}

\subsection{要旨}\label{ux8981ux65e8}

本稿では、先行論文 {[}1{]}, {[}2{]}, {[}3{]}, {[}4{]}, {[}5{]} で確立された中心投影フレームワークおよび次元追加の解釈的構造に基づき、ゼロ次元から四次元までの主観空間について、各次元で生まれる幾何学的性質と定説とされている構造との対応を記述する。

\textbf{明記事項}: 先行論文 {[}4{]}, {[}5{]} により、中心投影フレームワークにおける次元追加は原理的に無限に実行可能である。どこまで次元を追加するか、および追加の停止条件をどう解釈するかは解釈問題であり、本稿の考察範囲外である。本稿は考察の便宜上、ゼロ次元から四次元までの主観空間を扱うが、四次元という数は本稿の記述上の都合にすぎず、四次元で次元追加が停止すべき幾何学的理由を主張するものではない。

本稿では何も新たに証明しない。各次元の主観空間について、当該次元における定説とされる幾何学的性質が中心投影フレームワークの主観空間にどのように対応するかを記述するのみである。物理的解釈は本稿の対象外であり、読者に委ねる。

\begin{center}\rule{0.5\linewidth}{0.5pt}\end{center}

\subsection{1. 導入}\label{ux5c0eux5165}

先行論文 {[}1{]} で導入された中心投影フレームワークにおいて、主観空間 \(\Pi_R\) は曲率半径 \(R\) に支配される空間として定義される。論文 {[}2{]} は主観空間の軸の間の対等性を示し、論文 {[}3{]} は複数主観空間の観測の相対性を示した。論文 {[}4{]} は曲率半径の極限と入れ子構造の存在を示し、論文 {[}5{]} は次元追加の解釈的制約を示した。

本稿はこれらの結果に基づき、主観空間の次元を \(0, 1, 2, 3, 4\) と段階的に検討する。各次元について、当該次元で新たに生まれる幾何学的性質、認知可能になる量、および時間軸の解釈に関する観察を記述する。

繰り返しになるが、本稿はこれらの対応について何も新たに証明しない。各次元における記述は、当該次元の定説とされる幾何学的性質を主観空間に適用しただけのものである。本稿の意図は、中心投影フレームワークと定説との対応を整理することに限定される。

\begin{center}\rule{0.5\linewidth}{0.5pt}\end{center}

\subsection{2. ゼロ次元主観空間}\label{ux30bcux30edux6b21ux5143ux4e3bux89b3ux7a7aux9593}

ゼロ次元の主観空間は「ある / ない」のみを区別する世界である。

\textbf{内部構造}: 存在しない。

\textbf{認知可能な量}: 測地線偏差は存在しない。自身が曲率半径 \(R\) に支配されているかどうかを内部から認知する手段も存在しない。

\textbf{時間軸の解釈}: 該当しない(軸が存在しない)。

\textbf{次元追加の可能性}: 先行論文 {[}5{]} に従い、次元は自由に追加可能である。

\begin{center}\rule{0.5\linewidth}{0.5pt}\end{center}

\subsection{3. 一次元主観空間}\label{ux4e00ux6b21ux5143ux4e3bux89b3ux7a7aux9593}

一次元の主観空間は、\(0\) 以上の実数のみからなる世界である。

\textbf{内部構造}: 一本の軸。しかし基準点も方向の判断基準も存在しない。純粋な非負実数の集合である。長さの概念は存在するが、方向の概念は存在しない。

\textbf{認知可能な量}: 測地線偏差はまだ認知できない(一次元では偏差を定義する他の軸がない)。

\textbf{時間軸の解釈}: 該当しない(軸が一本しかなく、比較対象が存在しない)。

\textbf{次元追加の可能性}: 先行論文 {[}5{]} に従い、次元は自由に追加可能である。

\begin{center}\rule{0.5\linewidth}{0.5pt}\end{center}

\subsection{4. 二次元主観空間}\label{ux4e8cux6b21ux5143ux4e3bux89b3ux7a7aux9593}

二次元の主観空間で、面積と角度が新たに生まれる。

\textbf{内部構造}: 二本の軸。面が定義され、面に表裏が生まれる。この面の表裏により、一次元の長さに\textbf{方向}が付与される。ピタゴラスの定理が成立する。

\textbf{認知可能な量}: 面積、角度、および測地線偏差が認知可能になる。測地線偏差は二次元で初めて定義可能となり、論文 {[}2{]} の対応する命題が意味を持つ。

\textbf{時間軸の解釈}: 二本の軸は論文 {[}2{]} 定理 3.1 により対等であり、どちらを時間軸と解釈するかは決定できない。

\textbf{次元追加の可能性}: 先行論文 {[}5{]} に従い、次元は自由に追加可能である。

\begin{center}\rule{0.5\linewidth}{0.5pt}\end{center}

\subsection{5. 三次元主観空間}\label{ux4e09ux6b21ux5143ux4e3bux89b3ux7a7aux9593}

三次元の主観空間で、体積が新たに生まれる。

\textbf{内部構造}: 三本の軸。体積が定義され、面積に表裏が生まれる。

\textbf{認知可能な量}: 体積、面積、角度、測地線偏差が認知可能。測地線偏差は三次元でより豊かな構造を持つ。

\textbf{時間軸の解釈}: 三本の軸のうち一本を時間軸と解釈すると、残りの二本が空間軸として認知される。どの軸を時間軸と解釈するかは論文 {[}2{]} 定理 3.1 により観測者の選択に委ねられる。

\textbf{次元追加の可能性}: 先行論文 {[}5{]} に従い、次元は自由に追加可能である。

\begin{center}\rule{0.5\linewidth}{0.5pt}\end{center}

\subsection{6. 四次元主観空間}\label{ux56dbux6b21ux5143ux4e3bux89b3ux7a7aux9593}

四次元の主観空間で、四次元超体積が新たに生まれる。

\textbf{内部構造}: 四本の軸。四次元超体積が定義される。三次元体積に表裏が生まれる。

\textbf{認知可能な量}: 四次元超体積、三次元体積、面積、角度、測地線偏差が認知可能。

\textbf{時間軸の解釈}: 四本の軸のうち一本を時間軸と解釈すると、残りの三本が空間軸として認知される。一方、二本の軸を時間軸と解釈すると、残り二本の空間軸と幾何学的に区別できない(論文 {[}2{]} 定理 3.1 の軸の対等性による)。すなわち、時間軸として解釈される軸の本数は、四次元主観空間の内部では一意に決まらない。

\textbf{次元追加の可能性}: 先行論文 {[}5{]} に従い、次元は自由に追加可能である。

\begin{center}\rule{0.5\linewidth}{0.5pt}\end{center}

\subsection{7. 考察}\label{ux8003ux5bdf}

本稿では、ゼロ次元から四次元までの主観空間について、各次元で新たに生まれる幾何学的性質と定説との対応を記述した。

本稿で記述したのは以下に限定される。

\begin{enumerate}
\def\labelenumi{\arabic{enumi}.}
\tightlist
\item
  各次元で新たに定義可能になる量(面積、角度、体積、超体積、および対応する表裏構造)。
\item
  各次元で認知可能となる測地線偏差の状況。
\item
  各次元における時間軸の解釈可能性。
\item
  各段階で次元追加が自由であること。
\end{enumerate}

本稿で記述した各次元の性質は、当該次元における定説とされる幾何学的構造の主観空間への対応であり、新たな証明を含まない。これらの対応が持ちうる物理的意味付けについては、本稿では扱わない。特に、各次元で生まれる表裏構造が観測者の認知にどのように現れるか、時間軸の解釈がどのような基準で実現されるか、などの問いは幾何学の枠組みの外側の問題であり、本稿の射程を超える。

これらの問いへの解答は、読者の判断に委ねる。

\begin{center}\rule{0.5\linewidth}{0.5pt}\end{center}

\subsection{8. 結論}\label{ux7d50ux8ad6}

中心投影フレームワークの主観空間は、ゼロ次元から四次元まで、各次元において定説とされる幾何学的性質と自然に対応する。各段階で次元追加は自由であり、本稿の四次元までの記述は考察上の便宜にすぎない。各次元で現れる構造の物理的解釈は本稿の対象外であり、読者に委ねる。

\begin{center}\rule{0.5\linewidth}{0.5pt}\end{center}

\subsection{参考文献}\label{ux53c2ux8003ux6587ux732e}

{[}1{]} 木原 範昭,「中心投影による幾何学的定式化」, Zenodo, DOI: 10.5281/zenodo.19427780 (2025).

{[}2{]} 木原 範昭,「中心投影フレームワークにおける対称性」, Zenodo, DOI: 10.5281/zenodo.19434932 (2025).

{[}3{]} 木原 範昭,「複数主観空間における観測の相対性」, Zenodo, DOI: 10.5281/zenodo.19435162 (2025).

{[}4{]} 木原 範昭,「中心投影における曲率半径の極限についての考察 --- \(R \to \infty\) と \(R \to 0\) の場合」, Zenodo, DOI: 10.5281/zenodo.19526549 (2026).

{[}5{]} 木原 範昭,「背景空間と主観空間への次元追加についての考察」, Zenodo, DOI: 10.5281/zenodo.19526913 (2026).

\begin{center}\rule{0.5\linewidth}{0.5pt}\end{center}

\end{document}
